\chapter{Évaluation du Module 3}
\label{chap:evalmod3}

Cette évaluation valide les compétences liées à la gestion opérationnelle des incidents, l'investigation numérique (Forensics) et la planification de la continuité d'activité.

% ------------------------------------------------------------------------------
\section{Partie A : QCM (25 Questions)}
% ------------------------------------------------------------------------------

\textbf{Consigne :} Cochez la ou les bonnes réponses.

\begin{enumerate}
    \item Un événement qui compromet la sécurité de l'information est appelé :
    \begin{itemize}
        \item[A.] Un accident.
        \item[B.] Un incident.
        \item[C.] Une alerte.
        \item[D.] Un log.
    \end{itemize}

    \item La norme dédiée à la gestion des incidents de sécurité est :
    \begin{itemize}
        \item[A.] ISO 27001.
        \item[B.] ISO 27035.
        \item[C.] ISO 22301.
        \item[D.] ISO 27005.
    \end{itemize}

    \item Quel élément est le plus volatile (se perd le plus vite à l'extinction) ?
    \begin{itemize}
        \item[A.] Le disque dur.
        \item[B.] La mémoire RAM.
        \item[C.] Les sauvegardes sur bande.
        \item[D.] Les logs systèmes.
    \end{itemize}

    \item Le RTO (Recovery Time Objective) mesure :
    \begin{itemize}
        \item[A.] La quantité de données perdues.
        \item[B.] La durée cible pour rétablir le service.
        \item[C.] La fréquence des sauvegardes.
        \item[D.] Le temps de détection de l'incident.
    \end{itemize}

    \item Une "Chaîne de Custodie" (Chain of Custody) sert à :
    \begin{itemize}
        \item[A.] Verrouiller les portes du local serveur.
        \item[B.] Garantir l'intégrité et la traçabilité des preuves numériques.
        \item[C.] Connecter les serveurs entre eux.
        \item[D.] Envoyer des alertes par email.
    \end{itemize}

    \item Le RPO (Recovery Point Objective) mesure :
    \begin{itemize}
        \item[A.] Le temps de restoration.
        \item[B.] La perte de données acceptable (exprimée en temps).
        \item[C.] La distance du site de secours.
        \item[D.] Le coût de l'incident.
    \end{itemize}

    \item Un SIEM est un outil qui permet de :
    \begin{itemize}
        \item[A.] Chiffrer les disques.
        \item[B.] Centraliser les logs et corréler les événements.
        \item[C.] Gérer les mots de passe.
        \item[D.] Simuler des attaques.
    \end{itemize}

    \item La première action lors de la réaction à un incident est généralement :
    \begin{itemize}
        \item[A.] L'éradication.
        \item[B.] Le confinement.
        \item[C.] La communication.
        \item[D.] La réinstallation.
    \end{itemize}

    \item L'analyse d'impact sur l'activité (BIA) permet d'identifier :
    \begin{itemize}
        \item[A.] Les virus présents.
        \item[B.] Les fonctions vitales de l'entreprise.
        \item[C.] Les employés malveillants.
        \item[D.] Les coûts des logiciels.
    \end{itemize}

    \item Un "Faux Positif" est :
    \begin{itemize}
        \item[A.] Une alerte qui signale une menace qui n'existe pas.
        \item[B.] Un virus non détecté.
        \item[C.] Une erreur de l'administrateur.
        \item[D.] Une sauvegarde corrompue.
    \end{itemize}

    \item Le PCA (Plan de Continuité d'Activité) concerne :
    \begin{itemize}
        \item[A.] La réparation des serveurs uniquement.
        \item[B.] Le maintien des activités métier pendant la panne.
        \item[C.] La mise à jour des antivirus.
        \item[D.] Le paiement des rançons.
    \end{itemize}

    \item Un "Site Chaud" (Hot Site) est un site de secours :
    \begin{itemize}
        \item[A.] Sans matériel.
        \item[B.] Avec du matériel éteint.
        \item[C.] Avec du matériel allumé et synchronisé en temps réel.
        \item[D.] Situé dans un pays froid.
    \end{itemize}

    \item En investigation, un "Hash" sert à :
    \begin{itemize}
        \item[A.] Cacher les données.
        \item[B.] Vérifier l'intégrité d'un fichier ou d'une image disque.
        \item[C.] Accélérer le processeur.
        \item[D.] Effacer les preuves.
    \end{itemize}

    \item Le CSIRT est :
    \begin{itemize}
        \item[A.] Un logiciel de bureautique.
        \item[B.] L'équipe de réponse aux incidents de sécurité.
        \item[C.] Un type de pare-feu.
        \item[D.] Une norme juridique.
    \end{itemize}

    \item Le RETEX (Retour d'Expérience) se fait :
    \begin{itemize}
        \item[A.] Avant l'incident.
        \item[B.] Pendant l'incident.
        \item[C.] Après la résolution de l'incident.
        \item[D.] Jamais.
    \end{itemize}

    \item Un "Playbook" est :
    \begin{itemize}
        \item[A.] Un manuel de procédure de réponse à incident.
        \item[B.] Un outil de hacking.
        \item[C.] Un journal intime de l'admin.
        \item[D.] Un contrat d'assurance.
    \end{itemize}

    \item Selon le RGPD, une violation de données doit être notifiée sous :
    \begin{itemize}
        \item[A.] 24 heures.
        \item[B.] 72 heures.
        \item[C.] 1 semaine.
        \item[D.] 1 mois.
    \end{itemize}

    \item L'ordre de volatilité recommande de collecter en premier :
    \begin{itemize}
        \item[A.] Les fichiers sur le disque.
        \item[B.] Le contenu de la RAM.
        \item[C.] Les emails archivés.
        \item[D.] Les sauvegardes externes.
    \end{itemize}

    \item Le PRA (Plan de Reprise d'Activité) est :
    \begin{itemize}
        \item[A.] La version informatique du PCA.
        \item[B.] Un plan de marketing.
        \item[C.] Un plan de recrutement.
        \item[D.] Une procédure de paie.
    \end{itemize}

    \item Le "Confinement" consiste à :
    \begin{itemize}
        \item[A.] Supprimer le virus.
        \item[B.] Limiter la propagation de l'attaque.
        \item[C.] Restaurer les sauvegardes.
        \item[D.] Éteindre la lumière.
    \end{itemize}

    \item Une image disque est :
    \begin{itemize}
        \item[A.] Une photo du serveur.
        \item[B.] Une copie exacte bit-à-bit du support.
        \item[C.] Un raccourci Windows.
        \item[D.] Un fichier .jpg.
    \end{itemize}

    \item Le SOC est :
    \begin{itemize}
        \item[A.] Une équipe de surveillance 24/7.
        \item[B.] Un serveur email.
        \item[C.] Un protocole réseau.
        \item[D.] Un système d'exploitation.
    \end{itemize}

    \item La différence entre PCA et PRA est :
    \begin{itemize}
        \item[A.] PCA = Métier / PRA = Technique.
        \item[B.] PCA = Technique / PRA = Métier.
        \item[C.] Il n'y a pas de différence.
        \item[D.] PCA est pour les PME, PRA pour les grandes entreprises.
    \end{itemize}

    \item Un "Artéfact" en forensique est :
    \begin{itemize}
        \item[A.] Une pièce de musée.
        \item[B.] Une trace laissée par une activité numérique.
        \item[C.] Une erreur de code.
        \item[D.] Un objet physique.
    \end{itemize}

    \item En cas de Ransomware, la première mesure technique recommandée est souvent :
    \begin{itemize}
        \item[A.] Payer la rançon.
        \item[B.] Débrancher la machine du réseau.
        \item[C.] Redémarrer le serveur.
        \item[D.] Appeler la presse.
    \end{itemize}
\end{enumerate}

% ------------------------------------------------------------------------------
\section{Partie B : Questions Ouvertes (10 Questions)}
% ------------------------------------------------------------------------------

\textbf{Consigne :} Répondez de manière synthétique (5-10 lignes par question).

\begin{enumerate}
    \item Expliquez la différence entre un "Événement", un "Incident" et une "Crise".
    \item Pourquoi est-il déconseillé d'éteindre un serveur compromis immédiatement lors d'une investigation forensique ?
    \item Définissez le RTO et le RPO et expliquez leur relation avec le coût de la solution technique.
    \item Quelle est l'utilité d'un "Write Blocker" lors de l'acquisition d'une preuve numérique ?
    \item Expliquez le rôle du BIA dans la mise en place d'un PCA.
    \item Quelle est la différence entre un Site Froid, Tiède et Chaud ?
    \item Que signifie l'acronyme CSIRT et quel est son rôle principal ?
    \item Pourquoi le RETEX est-il une étape indispensable du cycle de vie de l'incident ?
    \item Expliquez la différence entre Correction et Éradication dans le cycle de réponse à incident.
    \item Quelles sont les informations minimales à inclure dans une fiche d'incident ?
\end{enumerate}

% ------------------------------------------------------------------------------
\section{Partie C : Cas Pratique - Mise en Condition (TP)}
% ------------------------------------------------------------------------------

\subsection{Scénario : Panne Critique du Datacenter}

\begin{boxlizbiz}
\textbf{Contexte :}
Un incendie s'est déclaré dans le local serveur principal de LiZbiZ pendant la nuit. Le système d'extinction a fonctionné, mais l'eau a endommagé l'ensemble des serveurs physiques. Le site Web est inaccessible, l'ERP est détruit. Les sauvegardes quotidiennes sont effectuées sur un Cloud externe sécurisé.
\end{boxlizbiz}

\subsection{Travaux à Réaliser (En temps limité)}

\textbf{Livrable 1 : Réaction Immédiate (30 min)}
En tant que DSI, vous coordonnez les premières heures.
\begin{itemize}
    \item Listez les 3 premières actions à réaliser (technique et humain).
    \item Définissez qui compose la Cellule de Crise (Rôles, pas de noms).
\end{itemize}

\textbf{Livrable 2 : Stratégie de Reprise (1h)}
Le PDG vous demande quand le site marchand sera de retour.
\begin{itemize}
    \item En vous basant sur le RTO défini (Chapitre 4), calculez l'impact si la restauration prend 10h.
    \item Proposez la procédure technique de restauration (PRA) depuis le Cloud.
    \item Proposez une solution de "Mode Dégradé" pour les commerciaux qui doivent répondre aux clients par téléphone.
\end{itemize}

\textbf{Livrable 3 : Rapport d'Incident (30 min)}
Rédigez la "Fiche d'Incident" finale (brève) qui servira au RETEX.
\begin{itemize}
    \item Date, Heure, Type d'incident.
    \item Cause racine (Rupture de tuyau ? Défaillance extincteur ?).
    \item Actions menées.
    \item Améliorations proposées (ex: Duplication du site).
\end{itemize}

% ==============================================================================
% Fichier : 05_module_crise.tex (Extrait Chapitre 1)
% Description : Introduction à la Gestion de Crise
% ==============================================================================